\documentclass[11pt,fleqn]{article}
\usepackage[utf8]{inputenc}
\usepackage[T1]{fontenc}
\usepackage{amsmath, amssymb, amsthm}
\usepackage{graphicx}
\usepackage{geometry}
\usepackage{xcolor}
\usepackage{hyperref}
\usepackage{booktabs}
\usepackage{caption}
\usepackage{fancyhdr}
\usepackage{float}

\geometry{a4paper, margin=1in}

\hypersetup{
    colorlinks=true,
    linkcolor=blue,
    citecolor=blue,
    urlcolor=blue
}

\pagestyle{fancy}
\fancyhf{}
\rhead{K-Dense AI Technical Report}
\lhead{Hamiltonians of Neural Topology}
\rfoot{Page \thepage}
\lfoot{Generated using K-Dense Web (\href{https://k-dense.ai}{k-dense.ai})}

\title{\textbf{The Hamiltonians of Neural Topology:}\\ \large A Symbolic Formalization of the $N_{\text{brain}}$ Tuple}
\author{K-Dense Web \\ \href{mailto:contact@k-dense.ai}{contact@k-dense.ai}}
\date{February 4, 2026}

\begin{document}

\maketitle

\begin{abstract}
This report distills the topological and dynamical principles of complex brain networks into a rigorous symbolic framework. We define the Neural Network System as a formal tuple $N_{\text{brain}} = \langle G, H, T, \Delta \rangle$, integrating graph geometry, evolutionary mechanics, functional scaling, and pathological perturbations. By establishing an optimization objective within an evolutionary Hamiltonian, we provide a symbolic proof for the emergence of the "Small-World" archetype as a critical state between lattice-like regularity and random stochasticity. Finally, we model the vulnerability of network hubs under targeted attacks, formalizing the dysconnectivity syndromes characteristic of neurodegenerative and psychiatric pathologies.
\end{abstract}

\begin{figure}[H]
    \centering
    \includegraphics[width=0.8\textwidth]{figures/graphical_abstract.png}
    \caption{\textbf{Graphical Abstract:} The $N_{\text{brain}}$ tuple spanning G (Geometry), H (Hamiltonian), T (Topology), and $\Delta$ (Perturbation). Figure generated using K-Dense Web.}
    \label{fig:graph_abstract}
\end{figure}

\section{Legend of Variables}
\begin{table}[H]
\centering
\begin{tabular}{ll}
\toprule
\textbf{Symbol} & \textbf{Definition} \\
\midrule
$G = (V, E)$ & Graph representation (Nodes $V$, Edges $E$) \\
$a_{ij} \in \{0, 1\}$ & Adjacency matrix element between nodes $i$ and $j$ \\
$d_{ij}$ & Physical Euclidean distance between nodes $i$ and $j$ \\
$P(a_{ij}=1)$ & Probability of connection as a function of distance \\
$C_{\text{wire}}$ & Wiring cost functional (Axonal Volume) \\
$L$ & Characteristic path length (average shortest path) \\
$E_{\text{glob}}$ & Global efficiency (inverse path length) \\
$H_{\text{evol}}$ & Evolutionary Hamiltonian (Cost-Efficiency trade-off) \\
$P(k)$ & Degree distribution (Truncated power-law) \\
$F_{ij}$ & Functional connectivity (Statistical correlation) \\
$\Psi_{\text{patho}}$ & Pathological operator (Dysconnectivity syndrome) \\
$\eta, \gamma, k_c$ & Constant parameters (Decay, exponent, cutoff) \\
\bottomrule
\end{tabular}
\end{table}

\section{Structural Probability Space ($G$)}
The structural substrate of the neural system is formalized as a graph $G = (V, E)$ embedded in a 3D physical coordinate space $\mathbb{R}^3$.

\subsection{Connection Probability}
A fundamental constraint on neural architecture is the energetic cost of long-range projections. The connection probability $P(a_{ij}=1)$ is modeled as a decaying function of physical distance $d_{ij}$:
\begin{equation}
    P(a_{ij} = 1) = \kappa \cdot e^{-\eta d_{ij}}
\end{equation}
where $\eta$ is the spatial decay constant, representing the metabolic penalty for distance.

\subsection{The Cost Functional}
The total wiring cost (axonal volume) $C_{\text{wire}}$ is defined as the sum of physical distances of all instantiated connections:
\begin{equation}
    C_{\text{wire}} = \sum_{i,j \in V} a_{ij} \cdot d_{ij}
\end{equation}
In biological systems, $C_{\text{wire}}$ is subject to strict minimization pressures, favoring local clustering and short-range connectivity [1].

\section{Evolutionary Hamiltonian ($H$)}
The configuration of $G$ is determined by an evolutionary objective function that balances topological efficiency against metabolic cost.

\subsection{Global Efficiency}
Topological efficiency $E_{\text{glob}}$ is defined as the average of the inverse shortest path lengths $L_{ij}$:
\begin{equation}
    E_{\text{glob}} = \frac{1}{N(N-1)} \sum_{i \neq j} \frac{1}{L_{ij}}
\end{equation}
Unlike path length $L$, $E_{\text{glob}}$ is well-defined even for disconnected graphs.

\subsection{The Economic Trade-off}
The Evolutionary Hamiltonian $H_{\text{evol}}$ quantifies the "fitness" of the network topology:
\begin{equation}
    \max H_{\text{evol}} = \alpha E_{\text{glob}} - \beta C_{\text{wire}}
\end{equation}
where $\alpha$ is the prize for communication efficiency and $\beta$ is the penalty for physical material cost.

\subsection{Proof of the Archetype Solution}
We examine the limit behaviors of $H_{\text{evol}}$:
\begin{enumerate}
    \item \textbf{Lattice Regime ($\beta \to \infty$):} Minimize $C_{\text{wire}}$ regardless of $E_{\text{glob}}$. The solution is a nearest-neighbor lattice where $C_{\text{wire}} \to \min$, but $E_{\text{glob}} \to 0$ as $N \to \infty$.
    \item \textbf{Random Regime ($\alpha \to \infty$):} Maximize $E_{\text{glob}}$ regardless of $C_{\text{wire}}$. The solution is a random graph where $L \sim \ln N$, resulting in high $E_{\text{glob}}$, but $C_{\text{wire}} \to \max$ due to long-range shortcuts.
\end{enumerate}
\textbf{The Small-World Archetype:} For finite $(\alpha, \beta)$, the system converges to a state with high local clustering (from $\beta$) and short path lengths (from $\alpha$). This "Small-World" topology satisfies the condition:
\begin{equation}
    \sigma = \frac{C_{\text{actual}} / C_{\text{random}}}{L_{\text{actual}} / L_{\text{random}}} > 1
\end{equation}
Thus, the brain network is a "critical" solution that optimizes information transfer at near-minimal cost [2].

\section{Functional Dynamics and Scaling ($T$)}
Functional dynamics represent the emergent properties of the structural substrate $G$.

\subsection{Connectivity Tensors}
Functional Connectivity $F_{ij}$ is defined as the time-averaged statistical correlation between states $x_i(t)$ and $x_j(t)$:
\begin{equation}
    F_{ij} = \frac{\langle x_i(t)x_j(t) \rangle - \langle x_i(t) \rangle \langle x_j(t) \rangle}{\sigma_i \sigma_j}
\end{equation}
We postulate a structure-function coupling $F_{ij} \approx f(A_{ij})$ at low frequencies. However, metastable dynamics allow for $dF/dt \neq 0$ while the hardware $dA/dt \approx 0$ on short timescales [4].

\subsection{The Truncated Power Law}
Unlike purely scale-free networks, physically embedded brain networks exhibit a degree distribution $P(k)$ following an exponentially truncated power law:
\begin{equation}
    P(k) \propto k^{-\gamma} e^{-k/k_c}
\end{equation}
where $k_c$ is the cutoff parameter. This truncation arises from physical constraints (spatial volume and metabolic saturation) that prevent the formation of infinite-capacity hubs [3].

\section{Pathological Operator ($\Delta$)}
Disorders of the brain are formalized as perturbations to the optimal network tuple.

\subsection{Hub Vulnerability}
Neural hubs $H = \{v \mid k_v \gg \langle k \rangle\}$ are defined as nodes with high degree and high betweenness centrality. These hubs are high-cost, high-efficiency bottlenecks.

\subsection{Targeted Attack Modeling}
We define the Pathological Operator $\Delta$ as a deletion of hub nodes $v \in H$. The impact on global efficiency is denoted as:
\begin{equation}
    \Psi_{\text{patho}} = E_{\text{glob}}(G) - E_{\text{glob}}(G - \{v \in H\})
\end{equation}
In Alzheimer's disease and Schizophrenia, $\Delta$ represents a targeted attack on the high-cost structural core. While brain networks are robust to random errors ($\delta v_{\text{rand}}$), they are critically vulnerable to the deletion of hubs, leading to a catastrophic collapse of $E_{\text{glob}}$ and the manifestation of "dysconnectivity syndromes" [5].

\section*{References}
\noindent [1] Bullmore, E., \& Sporns, O. (2009). Complex brain networks: graph theoretical analysis of structural and functional systems. \textit{Nature Reviews Neuroscience}, 10(3), 186-198. \\
\noindent [2] Watts, D. J., \& Strogatz, S. H. (1998). Collective dynamics of ‘small-world’ networks. \textit{Nature}, 393(6684), 442-445. \\
\noindent [3] Achard, S., et al. (2006). A resilient, low-frequency, small-world human brain functional network with highly connected association cortical hubs. \textit{Journal of Neuroscience}, 26(1), 63-72. \\
\noindent [4] Honey, C. J., et al. (2007). Network structure of cerebral cortex shapes functional connectivity on multiple time scales. \textit{PNAS}, 104(24), 10240-10245. \\
\noindent [5] Bassett, D. S., et al. (2006). Adaptive reconfiguration of fractal small-world human brain functional networks. \textit{PNAS}, 103(51), 19518-19523.

\end{document}
